\documentclass[11pt]{article}

    \usepackage[breakable]{tcolorbox}
    \usepackage{parskip} % Stop auto-indenting (to mimic markdown behaviour)
    

    % Basic figure setup, for now with no caption control since it's done
    % automatically by Pandoc (which extracts ![](path) syntax from Markdown).
    \usepackage{graphicx}
    % Keep aspect ratio if custom image width or height is specified
    \setkeys{Gin}{keepaspectratio}
    % Maintain compatibility with old templates. Remove in nbconvert 6.0
    \let\Oldincludegraphics\includegraphics
    % Ensure that by default, figures have no caption (until we provide a
    % proper Figure object with a Caption API and a way to capture that
    % in the conversion process - todo).
    \usepackage{caption}
    \DeclareCaptionFormat{nocaption}{}
    \captionsetup{format=nocaption,aboveskip=0pt,belowskip=0pt}

    \usepackage{float}
    \floatplacement{figure}{H} % forces figures to be placed at the correct location
    \usepackage{xcolor} % Allow colors to be defined
    \usepackage{enumerate} % Needed for markdown enumerations to work
    \usepackage{geometry} % Used to adjust the document margins
    \usepackage{amsmath} % Equations
    \usepackage{amssymb} % Equations
    \usepackage{textcomp} % defines textquotesingle
    % Hack from http://tex.stackexchange.com/a/47451/13684:
    \AtBeginDocument{%
        \def\PYZsq{\textquotesingle}% Upright quotes in Pygmentized code
    }
    \usepackage{upquote} % Upright quotes for verbatim code
    \usepackage{eurosym} % defines \euro

    \usepackage{iftex}
    \ifPDFTeX
        \usepackage[T1]{fontenc}
        \IfFileExists{alphabeta.sty}{
              \usepackage{alphabeta}
          }{
              \usepackage[mathletters]{ucs}
              \usepackage[utf8x]{inputenc}
          }
    \else
        \usepackage{fontspec}
        \usepackage{unicode-math}
    \fi

    \usepackage{fancyvrb} % verbatim replacement that allows latex
    \usepackage{grffile} % extends the file name processing of package graphics
                         % to support a larger range
    \makeatletter % fix for old versions of grffile with XeLaTeX
    \@ifpackagelater{grffile}{2019/11/01}
    {
      % Do nothing on new versions
    }
    {
      \def\Gread@@xetex#1{%
        \IfFileExists{"\Gin@base".bb}%
        {\Gread@eps{\Gin@base.bb}}%
        {\Gread@@xetex@aux#1}%
      }
    }
    \makeatother
    \usepackage[Export]{adjustbox} % Used to constrain images to a maximum size
    \adjustboxset{max size={0.9\linewidth}{0.9\paperheight}}

    % The hyperref package gives us a pdf with properly built
    % internal navigation ('pdf bookmarks' for the table of contents,
    % internal cross-reference links, web links for URLs, etc.)
    \usepackage{hyperref}
    % The default LaTeX title has an obnoxious amount of whitespace. By default,
    % titling removes some of it. It also provides customization options.
    \usepackage{titling}
    \usepackage{longtable} % longtable support required by pandoc >1.10
    \usepackage{booktabs}  % table support for pandoc > 1.12.2
    \usepackage{array}     % table support for pandoc >= 2.11.3
    \usepackage{calc}      % table minipage width calculation for pandoc >= 2.11.1
    \usepackage[inline]{enumitem} % IRkernel/repr support (it uses the enumerate* environment)
    \usepackage[normalem]{ulem} % ulem is needed to support strikethroughs (\sout)
                                % normalem makes italics be italics, not underlines
    \usepackage{soul}      % strikethrough (\st) support for pandoc >= 3.0.0
    \usepackage{mathrsfs}
    

    
    % Colors for the hyperref package
    \definecolor{urlcolor}{rgb}{0,.145,.698}
    \definecolor{linkcolor}{rgb}{.71,0.21,0.01}
    \definecolor{citecolor}{rgb}{.12,.54,.11}

    % ANSI colors
    \definecolor{ansi-black}{HTML}{3E424D}
    \definecolor{ansi-black-intense}{HTML}{282C36}
    \definecolor{ansi-red}{HTML}{E75C58}
    \definecolor{ansi-red-intense}{HTML}{B22B31}
    \definecolor{ansi-green}{HTML}{00A250}
    \definecolor{ansi-green-intense}{HTML}{007427}
    \definecolor{ansi-yellow}{HTML}{DDB62B}
    \definecolor{ansi-yellow-intense}{HTML}{B27D12}
    \definecolor{ansi-blue}{HTML}{208FFB}
    \definecolor{ansi-blue-intense}{HTML}{0065CA}
    \definecolor{ansi-magenta}{HTML}{D160C4}
    \definecolor{ansi-magenta-intense}{HTML}{A03196}
    \definecolor{ansi-cyan}{HTML}{60C6C8}
    \definecolor{ansi-cyan-intense}{HTML}{258F8F}
    \definecolor{ansi-white}{HTML}{C5C1B4}
    \definecolor{ansi-white-intense}{HTML}{A1A6B2}
    \definecolor{ansi-default-inverse-fg}{HTML}{FFFFFF}
    \definecolor{ansi-default-inverse-bg}{HTML}{000000}

    % common color for the border for error outputs.
    \definecolor{outerrorbackground}{HTML}{FFDFDF}

    % commands and environments needed by pandoc snippets
    % extracted from the output of `pandoc -s`
    \providecommand{\tightlist}{%
      \setlength{\itemsep}{0pt}\setlength{\parskip}{0pt}}
    \DefineVerbatimEnvironment{Highlighting}{Verbatim}{commandchars=\\\{\}}
    % Add ',fontsize=\small' for more characters per line
    \newenvironment{Shaded}{}{}
    \newcommand{\KeywordTok}[1]{\textcolor[rgb]{0.00,0.44,0.13}{\textbf{{#1}}}}
    \newcommand{\DataTypeTok}[1]{\textcolor[rgb]{0.56,0.13,0.00}{{#1}}}
    \newcommand{\DecValTok}[1]{\textcolor[rgb]{0.25,0.63,0.44}{{#1}}}
    \newcommand{\BaseNTok}[1]{\textcolor[rgb]{0.25,0.63,0.44}{{#1}}}
    \newcommand{\FloatTok}[1]{\textcolor[rgb]{0.25,0.63,0.44}{{#1}}}
    \newcommand{\CharTok}[1]{\textcolor[rgb]{0.25,0.44,0.63}{{#1}}}
    \newcommand{\StringTok}[1]{\textcolor[rgb]{0.25,0.44,0.63}{{#1}}}
    \newcommand{\CommentTok}[1]{\textcolor[rgb]{0.38,0.63,0.69}{\textit{{#1}}}}
    \newcommand{\OtherTok}[1]{\textcolor[rgb]{0.00,0.44,0.13}{{#1}}}
    \newcommand{\AlertTok}[1]{\textcolor[rgb]{1.00,0.00,0.00}{\textbf{{#1}}}}
    \newcommand{\FunctionTok}[1]{\textcolor[rgb]{0.02,0.16,0.49}{{#1}}}
    \newcommand{\RegionMarkerTok}[1]{{#1}}
    \newcommand{\ErrorTok}[1]{\textcolor[rgb]{1.00,0.00,0.00}{\textbf{{#1}}}}
    \newcommand{\NormalTok}[1]{{#1}}

    % Additional commands for more recent versions of Pandoc
    \newcommand{\ConstantTok}[1]{\textcolor[rgb]{0.53,0.00,0.00}{{#1}}}
    \newcommand{\SpecialCharTok}[1]{\textcolor[rgb]{0.25,0.44,0.63}{{#1}}}
    \newcommand{\VerbatimStringTok}[1]{\textcolor[rgb]{0.25,0.44,0.63}{{#1}}}
    \newcommand{\SpecialStringTok}[1]{\textcolor[rgb]{0.73,0.40,0.53}{{#1}}}
    \newcommand{\ImportTok}[1]{{#1}}
    \newcommand{\DocumentationTok}[1]{\textcolor[rgb]{0.73,0.13,0.13}{\textit{{#1}}}}
    \newcommand{\AnnotationTok}[1]{\textcolor[rgb]{0.38,0.63,0.69}{\textbf{\textit{{#1}}}}}
    \newcommand{\CommentVarTok}[1]{\textcolor[rgb]{0.38,0.63,0.69}{\textbf{\textit{{#1}}}}}
    \newcommand{\VariableTok}[1]{\textcolor[rgb]{0.10,0.09,0.49}{{#1}}}
    \newcommand{\ControlFlowTok}[1]{\textcolor[rgb]{0.00,0.44,0.13}{\textbf{{#1}}}}
    \newcommand{\OperatorTok}[1]{\textcolor[rgb]{0.40,0.40,0.40}{{#1}}}
    \newcommand{\BuiltInTok}[1]{{#1}}
    \newcommand{\ExtensionTok}[1]{{#1}}
    \newcommand{\PreprocessorTok}[1]{\textcolor[rgb]{0.74,0.48,0.00}{{#1}}}
    \newcommand{\AttributeTok}[1]{\textcolor[rgb]{0.49,0.56,0.16}{{#1}}}
    \newcommand{\InformationTok}[1]{\textcolor[rgb]{0.38,0.63,0.69}{\textbf{\textit{{#1}}}}}
    \newcommand{\WarningTok}[1]{\textcolor[rgb]{0.38,0.63,0.69}{\textbf{\textit{{#1}}}}}


    % Define a nice break command that doesn't care if a line doesn't already
    % exist.
    \def\br{\hspace*{\fill} \\* }
    % Math Jax compatibility definitions
    \def\gt{>}
    \def\lt{<}
    \let\Oldtex\TeX
    \let\Oldlatex\LaTeX
    \renewcommand{\TeX}{\textrm{\Oldtex}}
    \renewcommand{\LaTeX}{\textrm{\Oldlatex}}
    % Document parameters
    % Document title
    \title{Blood Spatter Analysis}
    
    
    
    
    
    
    
% Pygments definitions
\makeatletter
\def\PY@reset{\let\PY@it=\relax \let\PY@bf=\relax%
    \let\PY@ul=\relax \let\PY@tc=\relax%
    \let\PY@bc=\relax \let\PY@ff=\relax}
\def\PY@tok#1{\csname PY@tok@#1\endcsname}
\def\PY@toks#1+{\ifx\relax#1\empty\else%
    \PY@tok{#1}\expandafter\PY@toks\fi}
\def\PY@do#1{\PY@bc{\PY@tc{\PY@ul{%
    \PY@it{\PY@bf{\PY@ff{#1}}}}}}}
\def\PY#1#2{\PY@reset\PY@toks#1+\relax+\PY@do{#2}}

\@namedef{PY@tok@w}{\def\PY@tc##1{\textcolor[rgb]{0.73,0.73,0.73}{##1}}}
\@namedef{PY@tok@c}{\let\PY@it=\textit\def\PY@tc##1{\textcolor[rgb]{0.24,0.48,0.48}{##1}}}
\@namedef{PY@tok@cp}{\def\PY@tc##1{\textcolor[rgb]{0.61,0.40,0.00}{##1}}}
\@namedef{PY@tok@k}{\let\PY@bf=\textbf\def\PY@tc##1{\textcolor[rgb]{0.00,0.50,0.00}{##1}}}
\@namedef{PY@tok@kp}{\def\PY@tc##1{\textcolor[rgb]{0.00,0.50,0.00}{##1}}}
\@namedef{PY@tok@kt}{\def\PY@tc##1{\textcolor[rgb]{0.69,0.00,0.25}{##1}}}
\@namedef{PY@tok@o}{\def\PY@tc##1{\textcolor[rgb]{0.40,0.40,0.40}{##1}}}
\@namedef{PY@tok@ow}{\let\PY@bf=\textbf\def\PY@tc##1{\textcolor[rgb]{0.67,0.13,1.00}{##1}}}
\@namedef{PY@tok@nb}{\def\PY@tc##1{\textcolor[rgb]{0.00,0.50,0.00}{##1}}}
\@namedef{PY@tok@nf}{\def\PY@tc##1{\textcolor[rgb]{0.00,0.00,1.00}{##1}}}
\@namedef{PY@tok@nc}{\let\PY@bf=\textbf\def\PY@tc##1{\textcolor[rgb]{0.00,0.00,1.00}{##1}}}
\@namedef{PY@tok@nn}{\let\PY@bf=\textbf\def\PY@tc##1{\textcolor[rgb]{0.00,0.00,1.00}{##1}}}
\@namedef{PY@tok@ne}{\let\PY@bf=\textbf\def\PY@tc##1{\textcolor[rgb]{0.80,0.25,0.22}{##1}}}
\@namedef{PY@tok@nv}{\def\PY@tc##1{\textcolor[rgb]{0.10,0.09,0.49}{##1}}}
\@namedef{PY@tok@no}{\def\PY@tc##1{\textcolor[rgb]{0.53,0.00,0.00}{##1}}}
\@namedef{PY@tok@nl}{\def\PY@tc##1{\textcolor[rgb]{0.46,0.46,0.00}{##1}}}
\@namedef{PY@tok@ni}{\let\PY@bf=\textbf\def\PY@tc##1{\textcolor[rgb]{0.44,0.44,0.44}{##1}}}
\@namedef{PY@tok@na}{\def\PY@tc##1{\textcolor[rgb]{0.41,0.47,0.13}{##1}}}
\@namedef{PY@tok@nt}{\let\PY@bf=\textbf\def\PY@tc##1{\textcolor[rgb]{0.00,0.50,0.00}{##1}}}
\@namedef{PY@tok@nd}{\def\PY@tc##1{\textcolor[rgb]{0.67,0.13,1.00}{##1}}}
\@namedef{PY@tok@s}{\def\PY@tc##1{\textcolor[rgb]{0.73,0.13,0.13}{##1}}}
\@namedef{PY@tok@sd}{\let\PY@it=\textit\def\PY@tc##1{\textcolor[rgb]{0.73,0.13,0.13}{##1}}}
\@namedef{PY@tok@si}{\let\PY@bf=\textbf\def\PY@tc##1{\textcolor[rgb]{0.64,0.35,0.47}{##1}}}
\@namedef{PY@tok@se}{\let\PY@bf=\textbf\def\PY@tc##1{\textcolor[rgb]{0.67,0.36,0.12}{##1}}}
\@namedef{PY@tok@sr}{\def\PY@tc##1{\textcolor[rgb]{0.64,0.35,0.47}{##1}}}
\@namedef{PY@tok@ss}{\def\PY@tc##1{\textcolor[rgb]{0.10,0.09,0.49}{##1}}}
\@namedef{PY@tok@sx}{\def\PY@tc##1{\textcolor[rgb]{0.00,0.50,0.00}{##1}}}
\@namedef{PY@tok@m}{\def\PY@tc##1{\textcolor[rgb]{0.40,0.40,0.40}{##1}}}
\@namedef{PY@tok@gh}{\let\PY@bf=\textbf\def\PY@tc##1{\textcolor[rgb]{0.00,0.00,0.50}{##1}}}
\@namedef{PY@tok@gu}{\let\PY@bf=\textbf\def\PY@tc##1{\textcolor[rgb]{0.50,0.00,0.50}{##1}}}
\@namedef{PY@tok@gd}{\def\PY@tc##1{\textcolor[rgb]{0.63,0.00,0.00}{##1}}}
\@namedef{PY@tok@gi}{\def\PY@tc##1{\textcolor[rgb]{0.00,0.52,0.00}{##1}}}
\@namedef{PY@tok@gr}{\def\PY@tc##1{\textcolor[rgb]{0.89,0.00,0.00}{##1}}}
\@namedef{PY@tok@ge}{\let\PY@it=\textit}
\@namedef{PY@tok@gs}{\let\PY@bf=\textbf}
\@namedef{PY@tok@ges}{\let\PY@bf=\textbf\let\PY@it=\textit}
\@namedef{PY@tok@gp}{\let\PY@bf=\textbf\def\PY@tc##1{\textcolor[rgb]{0.00,0.00,0.50}{##1}}}
\@namedef{PY@tok@go}{\def\PY@tc##1{\textcolor[rgb]{0.44,0.44,0.44}{##1}}}
\@namedef{PY@tok@gt}{\def\PY@tc##1{\textcolor[rgb]{0.00,0.27,0.87}{##1}}}
\@namedef{PY@tok@err}{\def\PY@bc##1{{\setlength{\fboxsep}{\string -\fboxrule}\fcolorbox[rgb]{1.00,0.00,0.00}{1,1,1}{\strut ##1}}}}
\@namedef{PY@tok@kc}{\let\PY@bf=\textbf\def\PY@tc##1{\textcolor[rgb]{0.00,0.50,0.00}{##1}}}
\@namedef{PY@tok@kd}{\let\PY@bf=\textbf\def\PY@tc##1{\textcolor[rgb]{0.00,0.50,0.00}{##1}}}
\@namedef{PY@tok@kn}{\let\PY@bf=\textbf\def\PY@tc##1{\textcolor[rgb]{0.00,0.50,0.00}{##1}}}
\@namedef{PY@tok@kr}{\let\PY@bf=\textbf\def\PY@tc##1{\textcolor[rgb]{0.00,0.50,0.00}{##1}}}
\@namedef{PY@tok@bp}{\def\PY@tc##1{\textcolor[rgb]{0.00,0.50,0.00}{##1}}}
\@namedef{PY@tok@fm}{\def\PY@tc##1{\textcolor[rgb]{0.00,0.00,1.00}{##1}}}
\@namedef{PY@tok@vc}{\def\PY@tc##1{\textcolor[rgb]{0.10,0.09,0.49}{##1}}}
\@namedef{PY@tok@vg}{\def\PY@tc##1{\textcolor[rgb]{0.10,0.09,0.49}{##1}}}
\@namedef{PY@tok@vi}{\def\PY@tc##1{\textcolor[rgb]{0.10,0.09,0.49}{##1}}}
\@namedef{PY@tok@vm}{\def\PY@tc##1{\textcolor[rgb]{0.10,0.09,0.49}{##1}}}
\@namedef{PY@tok@sa}{\def\PY@tc##1{\textcolor[rgb]{0.73,0.13,0.13}{##1}}}
\@namedef{PY@tok@sb}{\def\PY@tc##1{\textcolor[rgb]{0.73,0.13,0.13}{##1}}}
\@namedef{PY@tok@sc}{\def\PY@tc##1{\textcolor[rgb]{0.73,0.13,0.13}{##1}}}
\@namedef{PY@tok@dl}{\def\PY@tc##1{\textcolor[rgb]{0.73,0.13,0.13}{##1}}}
\@namedef{PY@tok@s2}{\def\PY@tc##1{\textcolor[rgb]{0.73,0.13,0.13}{##1}}}
\@namedef{PY@tok@sh}{\def\PY@tc##1{\textcolor[rgb]{0.73,0.13,0.13}{##1}}}
\@namedef{PY@tok@s1}{\def\PY@tc##1{\textcolor[rgb]{0.73,0.13,0.13}{##1}}}
\@namedef{PY@tok@mb}{\def\PY@tc##1{\textcolor[rgb]{0.40,0.40,0.40}{##1}}}
\@namedef{PY@tok@mf}{\def\PY@tc##1{\textcolor[rgb]{0.40,0.40,0.40}{##1}}}
\@namedef{PY@tok@mh}{\def\PY@tc##1{\textcolor[rgb]{0.40,0.40,0.40}{##1}}}
\@namedef{PY@tok@mi}{\def\PY@tc##1{\textcolor[rgb]{0.40,0.40,0.40}{##1}}}
\@namedef{PY@tok@il}{\def\PY@tc##1{\textcolor[rgb]{0.40,0.40,0.40}{##1}}}
\@namedef{PY@tok@mo}{\def\PY@tc##1{\textcolor[rgb]{0.40,0.40,0.40}{##1}}}
\@namedef{PY@tok@ch}{\let\PY@it=\textit\def\PY@tc##1{\textcolor[rgb]{0.24,0.48,0.48}{##1}}}
\@namedef{PY@tok@cm}{\let\PY@it=\textit\def\PY@tc##1{\textcolor[rgb]{0.24,0.48,0.48}{##1}}}
\@namedef{PY@tok@cpf}{\let\PY@it=\textit\def\PY@tc##1{\textcolor[rgb]{0.24,0.48,0.48}{##1}}}
\@namedef{PY@tok@c1}{\let\PY@it=\textit\def\PY@tc##1{\textcolor[rgb]{0.24,0.48,0.48}{##1}}}
\@namedef{PY@tok@cs}{\let\PY@it=\textit\def\PY@tc##1{\textcolor[rgb]{0.24,0.48,0.48}{##1}}}

\def\PYZbs{\char`\\}
\def\PYZus{\char`\_}
\def\PYZob{\char`\{}
\def\PYZcb{\char`\}}
\def\PYZca{\char`\^}
\def\PYZam{\char`\&}
\def\PYZlt{\char`\<}
\def\PYZgt{\char`\>}
\def\PYZsh{\char`\#}
\def\PYZpc{\char`\%}
\def\PYZdl{\char`\$}
\def\PYZhy{\char`\-}
\def\PYZsq{\char`\'}
\def\PYZdq{\char`\"}
\def\PYZti{\char`\~}
% for compatibility with earlier versions
\def\PYZat{@}
\def\PYZlb{[}
\def\PYZrb{]}
\makeatother


    % For linebreaks inside Verbatim environment from package fancyvrb.
    \makeatletter
        \newbox\Wrappedcontinuationbox
        \newbox\Wrappedvisiblespacebox
        \newcommand*\Wrappedvisiblespace {\textcolor{red}{\textvisiblespace}}
        \newcommand*\Wrappedcontinuationsymbol {\textcolor{red}{\llap{\tiny$\m@th\hookrightarrow$}}}
        \newcommand*\Wrappedcontinuationindent {3ex }
        \newcommand*\Wrappedafterbreak {\kern\Wrappedcontinuationindent\copy\Wrappedcontinuationbox}
        % Take advantage of the already applied Pygments mark-up to insert
        % potential linebreaks for TeX processing.
        %        {, <, #, %, $, ' and ": go to next line.
        %        _, }, ^, &, >, - and ~: stay at end of broken line.
        % Use of \textquotesingle for straight quote.
        \newcommand*\Wrappedbreaksatspecials {%
            \def\PYGZus{\discretionary{\char`\_}{\Wrappedafterbreak}{\char`\_}}%
            \def\PYGZob{\discretionary{}{\Wrappedafterbreak\char`\{}{\char`\{}}%
            \def\PYGZcb{\discretionary{\char`\}}{\Wrappedafterbreak}{\char`\}}}%
            \def\PYGZca{\discretionary{\char`\^}{\Wrappedafterbreak}{\char`\^}}%
            \def\PYGZam{\discretionary{\char`\&}{\Wrappedafterbreak}{\char`\&}}%
            \def\PYGZlt{\discretionary{}{\Wrappedafterbreak\char`\<}{\char`\<}}%
            \def\PYGZgt{\discretionary{\char`\>}{\Wrappedafterbreak}{\char`\>}}%
            \def\PYGZsh{\discretionary{}{\Wrappedafterbreak\char`\#}{\char`\#}}%
            \def\PYGZpc{\discretionary{}{\Wrappedafterbreak\char`\%}{\char`\%}}%
            \def\PYGZdl{\discretionary{}{\Wrappedafterbreak\char`\$}{\char`\$}}%
            \def\PYGZhy{\discretionary{\char`\-}{\Wrappedafterbreak}{\char`\-}}%
            \def\PYGZsq{\discretionary{}{\Wrappedafterbreak\textquotesingle}{\textquotesingle}}%
            \def\PYGZdq{\discretionary{}{\Wrappedafterbreak\char`\"}{\char`\"}}%
            \def\PYGZti{\discretionary{\char`\~}{\Wrappedafterbreak}{\char`\~}}%
        }
        % Some characters . , ; ? ! / are not pygmentized.
        % This macro makes them "active" and they will insert potential linebreaks
        \newcommand*\Wrappedbreaksatpunct {%
            \lccode`\~`\.\lowercase{\def~}{\discretionary{\hbox{\char`\.}}{\Wrappedafterbreak}{\hbox{\char`\.}}}%
            \lccode`\~`\,\lowercase{\def~}{\discretionary{\hbox{\char`\,}}{\Wrappedafterbreak}{\hbox{\char`\,}}}%
            \lccode`\~`\;\lowercase{\def~}{\discretionary{\hbox{\char`\;}}{\Wrappedafterbreak}{\hbox{\char`\;}}}%
            \lccode`\~`\:\lowercase{\def~}{\discretionary{\hbox{\char`\:}}{\Wrappedafterbreak}{\hbox{\char`\:}}}%
            \lccode`\~`\?\lowercase{\def~}{\discretionary{\hbox{\char`\?}}{\Wrappedafterbreak}{\hbox{\char`\?}}}%
            \lccode`\~`\!\lowercase{\def~}{\discretionary{\hbox{\char`\!}}{\Wrappedafterbreak}{\hbox{\char`\!}}}%
            \lccode`\~`\/\lowercase{\def~}{\discretionary{\hbox{\char`\/}}{\Wrappedafterbreak}{\hbox{\char`\/}}}%
            \catcode`\.\active
            \catcode`\,\active
            \catcode`\;\active
            \catcode`\:\active
            \catcode`\?\active
            \catcode`\!\active
            \catcode`\/\active
            \lccode`\~`\~
        }
    \makeatother

    \let\OriginalVerbatim=\Verbatim
    \makeatletter
    \renewcommand{\Verbatim}[1][1]{%
        %\parskip\z@skip
        \sbox\Wrappedcontinuationbox {\Wrappedcontinuationsymbol}%
        \sbox\Wrappedvisiblespacebox {\FV@SetupFont\Wrappedvisiblespace}%
        \def\FancyVerbFormatLine ##1{\hsize\linewidth
            \vtop{\raggedright\hyphenpenalty\z@\exhyphenpenalty\z@
                \doublehyphendemerits\z@\finalhyphendemerits\z@
                \strut ##1\strut}%
        }%
        % If the linebreak is at a space, the latter will be displayed as visible
        % space at end of first line, and a continuation symbol starts next line.
        % Stretch/shrink are however usually zero for typewriter font.
        \def\FV@Space {%
            \nobreak\hskip\z@ plus\fontdimen3\font minus\fontdimen4\font
            \discretionary{\copy\Wrappedvisiblespacebox}{\Wrappedafterbreak}
            {\kern\fontdimen2\font}%
        }%

        % Allow breaks at special characters using \PYG... macros.
        \Wrappedbreaksatspecials
        % Breaks at punctuation characters . , ; ? ! and / need catcode=\active
        \OriginalVerbatim[#1,codes*=\Wrappedbreaksatpunct]%
    }
    \makeatother

    % Exact colors from NB
    \definecolor{incolor}{HTML}{303F9F}
    \definecolor{outcolor}{HTML}{D84315}
    \definecolor{cellborder}{HTML}{CFCFCF}
    \definecolor{cellbackground}{HTML}{F7F7F7}

    % prompt
    \makeatletter
    \newcommand{\boxspacing}{\kern\kvtcb@left@rule\kern\kvtcb@boxsep}
    \makeatother
    \newcommand{\prompt}[4]{
        {\ttfamily\llap{{\color{#2}[#3]:\hspace{3pt}#4}}\vspace{-\baselineskip}}
    }
    

    
    % Prevent overflowing lines due to hard-to-break entities
    \sloppy
    % Setup hyperref package
    \hypersetup{
      breaklinks=true,  % so long urls are correctly broken across lines
      colorlinks=true,
      urlcolor=urlcolor,
      linkcolor=linkcolor,
      citecolor=citecolor,
      }
    % Slightly bigger margins than the latex defaults
    
    \geometry{verbose,tmargin=1in,bmargin=1in,lmargin=1in,rmargin=1in}
    
    

\begin{document}
    
    \maketitle
    
    

    
  
\section{Content}\label{content}

\begin{enumerate}
\def\labelenumi{\arabic{enumi}.}
\tightlist
\item
  Importance of Blood Spatter Analysis
\item
  The Math Behind Blood Spatter Analysis

  \begin{itemize}
  \tightlist
  \item
    Trigonometry of Blood Droplets
  \item
    Determining the Point of Origin
  \item
    Blood Drop Velocity \& Force Estimation
  \item
    Blood Spatter Distribution \& Area of Convergence
  \end{itemize}
\item
  Data simulation
\item
  Extracting insights from blood spatter data

  \begin{itemize}
  \tightlist
  \item
    Reverse engineering (finding source of spatter)
  \end{itemize}
\item
  Machine learning
\item
  Assumptions of Blood Spatter Analysis
\item
  Conclusion
\end{enumerate}

\section{Importance of Blood Spatter
Analysis}\label{importance-of-blood-spatter-analysis}

Bloodstain pattern analysis (BPA) is crucial in forensic investigations
because it helps reconstruct the events of a crime. By analyzing
bloodstain patterns, forensic experts can determine: 1. The type of
force used (blunt force, sharp force, gunshot) 2. The position of the
victim and perpetrator 3. The movement of individuals after bloodshed 4.
The minimum number of blows or shots inflicted 5. Whether a crime scene
has been altered

BPA relies heavily on fluid dynamics, trigonometry, and physics, which
we will now explore in detail.

\section{The Math Behind Blood Spatter
Analysis}\label{the-math-behind-blood-spatter-analysis}

\subsection{A. Trigonometry of Blood
Droplets}\label{a.-trigonometry-of-blood-droplets}

When blood is projected, it follows the laws of physics. One of the key
aspects analyzed is the angle of impact (θ), which determines the
direction and origin of the bloodstain.

The angle of impact is calculated using the formula:

\[\theta = arcsin\left(\frac{width}{length}\right)\]

Where: - \emph{\(\theta\)} = Angle of impact (in degrees) - \emph{w} =
Width of the bloodstain - \emph{l} = Length of the bloodstain

\begin{itemize}
\tightlist
\item
  \emph{Key Insight}: The more elongated the stain, the smaller the
  angle of impact. A circular stain (where w=l) means the blood hit the
  surface at \textbf{90° (perpendicular impact)}.
\end{itemize}

\subsection{B. Determining the Point of
Origin}\label{b.-determining-the-point-of-origin}

To determine where the blood originated in 3D space, forensic experts
use triangulation.

\textbf{Height of Origin (Z-axis Calculation)} Once we have multiple
impact angles, we can estimate the height (h) from which blood
originated using the tangent function:

\[h=d.\tan(\theta)\]

Where: - \emph{h} = Height of origin - \emph{d} = Horizontal distance
from bloodstain to the assumed point of origin - \emph{\(\theta\)} =
Angle of impact

By analyzing multiple stains from different angles, we can pinpoint the
exact location of the blood source.

\subsection{C. Blood Drop Velocity \& Force
Estimation}\label{c.-blood-drop-velocity-force-estimation}

\begin{enumerate}
\def\labelenumi{\arabic{enumi}.}
\tightlist
\item
  \textbf{Terminal Velocity of Blood Drops}
\end{enumerate}

The velocity of a free-falling blood droplet follows standard physics
equations:

\[v=\sqrt{2gh}\]

Where:

\begin{itemize}
\item
  \emph{v} = Velocity of blood droplet
\item
  \emph{g} = Acceleration due to gravity (9.81 m/s²)
\item
  \emph{h} = Height from which the blood originated
\item
  \emph{Key Insight}: A droplet falling from a greater height will have
  a higher velocity, affecting stain size and distribution.
\end{itemize}

\begin{enumerate}
\def\labelenumi{\arabic{enumi}.}
\setcounter{enumi}{1}
\tightlist
\item
  \textbf{Estimating the Force Behind Blood Spatter} Different forces
  cause different types of spatter:
\end{enumerate}

\begin{itemize}
\tightlist
\item
  \textbf{Low-velocity spatter (\textless5~ft/s)}: Large stains (4--6 mm
  diameter), typically from blunt objects or passive bleeding.
\item
  \textbf{Medium-velocity spatter (5 to 25 ft/s)}: Smaller stains (1--4
  mm), often from beatings or stabbings.
\item
  \textbf{High-velocity spatter (\textgreater100~ft/s)}: Tiny mist-like
  droplets (\textless1~mm), usually from gunshots.
\end{itemize}

The kinetic energy of a moving weapon or projectile can be estimated
using:

\[KE=\frac{1}{2}mv^2\]

Where: - \emph{KE} = Kinetic Energy - \emph{m} = Mass of the
weapon/projectile - \emph{v} = Velocity of impact

\begin{itemize}
\tightlist
\item
  \emph{Key Insight}: Gunshots produce \textbf{high-velocity mist-like
  spatter}, while blunt force trauma leads to \textbf{medium-velocity
  spatter} with larger drops.
\end{itemize}

\subsection{D. Blood Spatter Distribution \& Area of
Convergence}\label{d.-blood-spatter-distribution-area-of-convergence}

If multiple bloodstains originate from the same source, forensic
analysts determine their point of convergence, which is the 2D location
from where they came.

Using vector analysis, each bloodstain's trajectory is drawn backward to
find the intersection:

\[X=\frac{\sum(d_i .\tan(\theta_i))}{n}\]

Where:

\begin{itemize}
\item
  \emph{X} = X-coordinate of the point of convergence
\item
  \emph{\(d_i\)} = Distance of each bloodstain from a reference point
\item
  \emph{\(θ_i\)} = Angle of impact of each bloodstain
\item
  \emph{n} = Number of bloodstains analyzed
\item
  \emph{Key Insight}: The point of convergence gives a 2D location,
  while the point of origin gives the 3D position.
\end{itemize}

    \section{Data Simulation}\label{data-simulation}

Using Python, we shall simulate blood spatter data. We shall generate
random blood droplets with \textbf{velocity}, \textbf{impact angle} and
\textbf{direction}. We shall follow by simulating blood spatter patterns
by calculating stain size, shape and distribution. Then we shall
visualize the spatter using Matplotlib.

    \begin{tcolorbox}[breakable, size=fbox, boxrule=1pt, pad at break*=1mm,colback=cellbackground, colframe=cellborder]
\prompt{In}{incolor}{2}{\boxspacing}
\begin{Verbatim}[commandchars=\\\{\}]
\PY{k+kn}{import}\PY{+w}{ }\PY{n+nn}{numpy}\PY{+w}{ }\PY{k}{as}\PY{+w}{ }\PY{n+nn}{np}
\PY{k+kn}{import}\PY{+w}{ }\PY{n+nn}{matplotlib}\PY{n+nn}{.}\PY{n+nn}{pyplot}\PY{+w}{ }\PY{k}{as}\PY{+w}{ }\PY{n+nn}{plt}

\PY{k}{def}\PY{+w}{ }\PY{n+nf}{generate\PYZus{}blood\PYZus{}spatter}\PY{p}{(}\PY{n}{num\PYZus{}drops}\PY{o}{=}\PY{l+m+mi}{100}\PY{p}{,} \PY{n}{origin}\PY{o}{=}\PY{p}{(}\PY{l+m+mi}{0}\PY{p}{,} \PY{l+m+mi}{0}\PY{p}{)}\PY{p}{,} \PY{n}{max\PYZus{}velocity}\PY{o}{=}\PY{l+m+mi}{10}\PY{p}{,} \PY{n}{spread}\PY{o}{=}\PY{l+m+mi}{30}\PY{p}{)}\PY{p}{:}
\PY{+w}{    }\PY{l+s+sd}{\PYZdq{}\PYZdq{}\PYZdq{}Simulates blood spatter using random angles and velocities.\PYZdq{}\PYZdq{}\PYZdq{}}
    \PY{n}{x\PYZus{}points}\PY{p}{,} \PY{n}{y\PYZus{}points} \PY{o}{=} \PY{p}{[}\PY{p}{]}\PY{p}{,} \PY{p}{[}\PY{p}{]}
    
    \PY{k}{for} \PY{n}{\PYZus{}} \PY{o+ow}{in} \PY{n+nb}{range}\PY{p}{(}\PY{n}{num\PYZus{}drops}\PY{p}{)}\PY{p}{:}
        \PY{c+c1}{\PYZsh{} Random velocity}
        \PY{n}{velocity} \PY{o}{=} \PY{n}{np}\PY{o}{.}\PY{n}{random}\PY{o}{.}\PY{n}{uniform}\PY{p}{(}\PY{l+m+mi}{1}\PY{p}{,} \PY{n}{max\PYZus{}velocity}\PY{p}{)}
        
        \PY{c+c1}{\PYZsh{} Random impact angle (in degrees, then converted to radians)}
        \PY{n}{angle} \PY{o}{=} \PY{n}{np}\PY{o}{.}\PY{n}{random}\PY{o}{.}\PY{n}{uniform}\PY{p}{(}\PY{o}{\PYZhy{}}\PY{n}{spread}\PY{p}{,} \PY{n}{spread}\PY{p}{)} \PY{o}{*} \PY{p}{(}\PY{n}{np}\PY{o}{.}\PY{n}{pi} \PY{o}{/} \PY{l+m+mi}{180}\PY{p}{)}
        
        \PY{c+c1}{\PYZsh{} Compute x, y displacement (assuming flat surface impact)}
        \PY{n}{x} \PY{o}{=} \PY{n}{origin}\PY{p}{[}\PY{l+m+mi}{0}\PY{p}{]} \PY{o}{+} \PY{n}{velocity} \PY{o}{*} \PY{n}{np}\PY{o}{.}\PY{n}{cos}\PY{p}{(}\PY{n}{angle}\PY{p}{)}
        \PY{n}{y} \PY{o}{=} \PY{n}{origin}\PY{p}{[}\PY{l+m+mi}{1}\PY{p}{]} \PY{o}{+} \PY{n}{velocity} \PY{o}{*} \PY{n}{np}\PY{o}{.}\PY{n}{sin}\PY{p}{(}\PY{n}{angle}\PY{p}{)}
        
        \PY{n}{x\PYZus{}points}\PY{o}{.}\PY{n}{append}\PY{p}{(}\PY{n}{x}\PY{p}{)}
        \PY{n}{y\PYZus{}points}\PY{o}{.}\PY{n}{append}\PY{p}{(}\PY{n}{y}\PY{p}{)}
    
    \PY{k}{return} \PY{n}{x\PYZus{}points}\PY{p}{,} \PY{n}{y\PYZus{}points}

\PY{c+c1}{\PYZsh{} Generate and plot blood spatter}
\PY{n}{x}\PY{p}{,} \PY{n}{y} \PY{o}{=} \PY{n}{generate\PYZus{}blood\PYZus{}spatter}\PY{p}{(}\PY{n}{num\PYZus{}drops}\PY{o}{=}\PY{l+m+mi}{200}\PY{p}{,} \PY{n}{max\PYZus{}velocity}\PY{o}{=}\PY{l+m+mi}{15}\PY{p}{,} \PY{n}{spread}\PY{o}{=}\PY{l+m+mi}{45}\PY{p}{)}
\PY{n}{plt}\PY{o}{.}\PY{n}{figure}\PY{p}{(}\PY{n}{figsize}\PY{o}{=}\PY{p}{(}\PY{l+m+mi}{9}\PY{p}{,} \PY{l+m+mi}{6}\PY{p}{)}\PY{p}{)}
\PY{n}{plt}\PY{o}{.}\PY{n}{scatter}\PY{p}{(}\PY{n}{x}\PY{p}{,} \PY{n}{y}\PY{p}{,} \PY{n}{color}\PY{o}{=}\PY{l+s+s1}{\PYZsq{}}\PY{l+s+s1}{red}\PY{l+s+s1}{\PYZsq{}}\PY{p}{,} \PY{n}{alpha}\PY{o}{=}\PY{l+m+mf}{0.6}\PY{p}{,} \PY{n}{s}\PY{o}{=}\PY{l+m+mi}{10}\PY{p}{)}
\PY{n}{plt}\PY{o}{.}\PY{n}{xlabel}\PY{p}{(}\PY{l+s+s2}{\PYZdq{}}\PY{l+s+s2}{X Position}\PY{l+s+s2}{\PYZdq{}}\PY{p}{)}
\PY{n}{plt}\PY{o}{.}\PY{n}{ylabel}\PY{p}{(}\PY{l+s+s2}{\PYZdq{}}\PY{l+s+s2}{Y Position}\PY{l+s+s2}{\PYZdq{}}\PY{p}{)}
\PY{n}{plt}\PY{o}{.}\PY{n}{title}\PY{p}{(}\PY{l+s+s2}{\PYZdq{}}\PY{l+s+s2}{Simulated Blood Spatter}\PY{l+s+s2}{\PYZdq{}}\PY{p}{)}
\PY{n}{plt}\PY{o}{.}\PY{n}{show}\PY{p}{(}\PY{p}{)}
\end{Verbatim}
\end{tcolorbox}

    \begin{center}
    \adjustimage{max size={0.9\linewidth}{0.9\paperheight}}{output_2_0.png}
    \end{center}
    { \hspace*{\fill} \\}
    
    \subsection{What type of data did we just
simulate?}\label{what-type-of-data-did-we-just-simulate}

The simulation involved: - A single origin point → Suggests a localized
source, which could be a wound. - Randomized velocities (1--15 m/s) →
Simulates different impact forces. - Random spread angles
(\(\pm45^{\circ}\)) → Creates a medium-width distribution.

Looking at these parameters, the generated pattern resembles: * Impact
Spatter rather than passive dripping. * Medium-velocity spatter (based
on velocity range). * A possible blunt force or stabbing rather than a
high-energy gunshot.

If we wanted to simulate data for, say a gunshot wound, we can modify
the code by increasing velocity range (30-100 m/s), reduce droplet size
and narrow the spread angle (\(\pm15^{\circ}\))

\subsection{Extracting Insights from Blood
Spatter}\label{extracting-insights-from-blood-spatter}

We need to analyze and understand the simulated spatter data.

The key analyses in this step involve: 1. \textbf{Droplet Size
Distribution} -- How large are the stains? (Histogram, mean, variance)
2. \textbf{Velocity vs.~Distance Relationship} -- Do faster droplets
travel farther? (Scatter plot, correlation) 3. \textbf{Impact Angle
Estimation} -- Use trigonometry to estimate the angle of impact 4.
\textbf{Find the Source of Blood (Reverse Engineering)} -- Use
trigonometry \& geometry

Using python, we generate a code to analyze the first three then proceed
to reverse engineering.

    \begin{tcolorbox}[breakable, size=fbox, boxrule=1pt, pad at break*=1mm,colback=cellbackground, colframe=cellborder]
\prompt{In}{incolor}{3}{\boxspacing}
\begin{Verbatim}[commandchars=\\\{\}]
\PY{k+kn}{import}\PY{+w}{ }\PY{n+nn}{numpy}\PY{+w}{ }\PY{k}{as}\PY{+w}{ }\PY{n+nn}{np}
\PY{k+kn}{import}\PY{+w}{ }\PY{n+nn}{matplotlib}\PY{n+nn}{.}\PY{n+nn}{pyplot}\PY{+w}{ }\PY{k}{as}\PY{+w}{ }\PY{n+nn}{plt}
\PY{k+kn}{import}\PY{+w}{ }\PY{n+nn}{seaborn}\PY{+w}{ }\PY{k}{as}\PY{+w}{ }\PY{n+nn}{sns}

\PY{c+c1}{\PYZsh{} Generate synthetic blood spatter data}
\PY{n}{np}\PY{o}{.}\PY{n}{random}\PY{o}{.}\PY{n}{seed}\PY{p}{(}\PY{l+m+mi}{42}\PY{p}{)}
\PY{n}{num\PYZus{}drops} \PY{o}{=} \PY{l+m+mi}{200}

\PY{n}{drop\PYZus{}sizes} \PY{o}{=} \PY{n}{np}\PY{o}{.}\PY{n}{random}\PY{o}{.}\PY{n}{normal}\PY{p}{(}\PY{n}{loc}\PY{o}{=}\PY{l+m+mi}{3}\PY{p}{,} \PY{n}{scale}\PY{o}{=}\PY{l+m+mi}{1}\PY{p}{,} \PY{n}{size}\PY{o}{=}\PY{n}{num\PYZus{}drops}\PY{p}{)}  \PY{c+c1}{\PYZsh{} Mean 3mm, Std 1mm}
\PY{n}{velocities} \PY{o}{=} \PY{n}{np}\PY{o}{.}\PY{n}{random}\PY{o}{.}\PY{n}{uniform}\PY{p}{(}\PY{l+m+mi}{2}\PY{p}{,} \PY{l+m+mi}{20}\PY{p}{,} \PY{n}{num\PYZus{}drops}\PY{p}{)}  \PY{c+c1}{\PYZsh{} Velocity range from 2\PYZhy{}20 m/s}
\PY{n}{distances} \PY{o}{=} \PY{n}{velocities} \PY{o}{*} \PY{n}{np}\PY{o}{.}\PY{n}{random}\PY{o}{.}\PY{n}{uniform}\PY{p}{(}\PY{l+m+mf}{0.5}\PY{p}{,} \PY{l+m+mf}{1.5}\PY{p}{,} \PY{n}{num\PYZus{}drops}\PY{p}{)}  \PY{c+c1}{\PYZsh{} Distance depends on velocity}
\PY{n}{angles} \PY{o}{=} \PY{n}{np}\PY{o}{.}\PY{n}{arcsin}\PY{p}{(}\PY{n}{drop\PYZus{}sizes} \PY{o}{/} \PY{p}{(}\PY{n}{drop\PYZus{}sizes} \PY{o}{+} \PY{n}{np}\PY{o}{.}\PY{n}{random}\PY{o}{.}\PY{n}{uniform}\PY{p}{(}\PY{l+m+mi}{1}\PY{p}{,} \PY{l+m+mi}{5}\PY{p}{,} \PY{n}{num\PYZus{}drops}\PY{p}{)}\PY{p}{)}\PY{p}{)} \PY{o}{*} \PY{p}{(}\PY{l+m+mi}{180} \PY{o}{/} \PY{n}{np}\PY{o}{.}\PY{n}{pi}\PY{p}{)}  \PY{c+c1}{\PYZsh{} Impact angles}

\PY{c+c1}{\PYZsh{} \PYZhy{}\PYZhy{}\PYZhy{} Plot 1: Droplet Size Distribution \PYZhy{}\PYZhy{}\PYZhy{}}
\PY{n}{sns}\PY{o}{.}\PY{n}{histplot}\PY{p}{(}\PY{n}{drop\PYZus{}sizes}\PY{p}{,} \PY{n}{bins}\PY{o}{=}\PY{l+m+mi}{15}\PY{p}{,} \PY{n}{kde}\PY{o}{=}\PY{k+kc}{True}\PY{p}{,} \PY{n}{color}\PY{o}{=}\PY{l+s+s1}{\PYZsq{}}\PY{l+s+s1}{red}\PY{l+s+s1}{\PYZsq{}}\PY{p}{)}
\PY{n}{plt}\PY{o}{.}\PY{n}{xlabel}\PY{p}{(}\PY{l+s+s1}{\PYZsq{}}\PY{l+s+s1}{Droplet Size (mm)}\PY{l+s+s1}{\PYZsq{}}\PY{p}{)}
\PY{n}{plt}\PY{o}{.}\PY{n}{ylabel}\PY{p}{(}\PY{l+s+s1}{\PYZsq{}}\PY{l+s+s1}{Frequency}\PY{l+s+s1}{\PYZsq{}}\PY{p}{)}
\PY{n}{plt}\PY{o}{.}\PY{n}{title}\PY{p}{(}\PY{l+s+s1}{\PYZsq{}}\PY{l+s+s1}{Droplet Size Distribution}\PY{l+s+s1}{\PYZsq{}}\PY{p}{)}
\PY{n}{plt}\PY{o}{.}\PY{n}{show}\PY{p}{(}\PY{p}{)}
\end{Verbatim}
\end{tcolorbox}

    \begin{center}
    \adjustimage{max size={0.9\linewidth}{0.9\paperheight}}{output_4_0.png}
    \end{center}
    { \hspace*{\fill} \\}
    
    \begin{tcolorbox}[breakable, size=fbox, boxrule=1pt, pad at break*=1mm,colback=cellbackground, colframe=cellborder]
\prompt{In}{incolor}{4}{\boxspacing}
\begin{Verbatim}[commandchars=\\\{\}]
\PY{c+c1}{\PYZsh{} \PYZhy{}\PYZhy{}\PYZhy{} Plot 2: Velocity vs. Distance Relationship \PYZhy{}\PYZhy{}\PYZhy{}}
\PY{n}{plt}\PY{o}{.}\PY{n}{scatter}\PY{p}{(}\PY{n}{velocities}\PY{p}{,} \PY{n}{distances}\PY{p}{,} \PY{n}{alpha}\PY{o}{=}\PY{l+m+mf}{0.7}\PY{p}{,} \PY{n}{color}\PY{o}{=}\PY{l+s+s1}{\PYZsq{}}\PY{l+s+s1}{blue}\PY{l+s+s1}{\PYZsq{}}\PY{p}{)}
\PY{n}{plt}\PY{o}{.}\PY{n}{xlabel}\PY{p}{(}\PY{l+s+s1}{\PYZsq{}}\PY{l+s+s1}{Velocity (m/s)}\PY{l+s+s1}{\PYZsq{}}\PY{p}{)}
\PY{n}{plt}\PY{o}{.}\PY{n}{ylabel}\PY{p}{(}\PY{l+s+s1}{\PYZsq{}}\PY{l+s+s1}{Distance (cm)}\PY{l+s+s1}{\PYZsq{}}\PY{p}{)}
\PY{n}{plt}\PY{o}{.}\PY{n}{title}\PY{p}{(}\PY{l+s+s1}{\PYZsq{}}\PY{l+s+s1}{Velocity vs. Distance Relationship}\PY{l+s+s1}{\PYZsq{}}\PY{p}{)}
\PY{n}{plt}\PY{o}{.}\PY{n}{show}\PY{p}{(}\PY{p}{)}
\end{Verbatim}
\end{tcolorbox}

    \begin{center}
    \adjustimage{max size={0.9\linewidth}{0.9\paperheight}}{output_5_0.png}
    \end{center}
    { \hspace*{\fill} \\}
    
    \begin{tcolorbox}[breakable, size=fbox, boxrule=1pt, pad at break*=1mm,colback=cellbackground, colframe=cellborder]
\prompt{In}{incolor}{5}{\boxspacing}
\begin{Verbatim}[commandchars=\\\{\}]
\PY{c+c1}{\PYZsh{} \PYZhy{}\PYZhy{}\PYZhy{} Plot 3: Impact Angle Estimation \PYZhy{}\PYZhy{}\PYZhy{}}
\PY{n}{sns}\PY{o}{.}\PY{n}{histplot}\PY{p}{(}\PY{n}{angles}\PY{p}{,} \PY{n}{bins}\PY{o}{=}\PY{l+m+mi}{15}\PY{p}{,} \PY{n}{kde}\PY{o}{=}\PY{k+kc}{True}\PY{p}{,} \PY{n}{color}\PY{o}{=}\PY{l+s+s1}{\PYZsq{}}\PY{l+s+s1}{green}\PY{l+s+s1}{\PYZsq{}}\PY{p}{)}
\PY{n}{plt}\PY{o}{.}\PY{n}{xlabel}\PY{p}{(}\PY{l+s+s1}{\PYZsq{}}\PY{l+s+s1}{Impact Angle (Degrees)}\PY{l+s+s1}{\PYZsq{}}\PY{p}{)}
\PY{n}{plt}\PY{o}{.}\PY{n}{ylabel}\PY{p}{(}\PY{l+s+s1}{\PYZsq{}}\PY{l+s+s1}{Frequency}\PY{l+s+s1}{\PYZsq{}}\PY{p}{)}
\PY{n}{plt}\PY{o}{.}\PY{n}{title}\PY{p}{(}\PY{l+s+s1}{\PYZsq{}}\PY{l+s+s1}{Impact Angle Distribution}\PY{l+s+s1}{\PYZsq{}}\PY{p}{)}
\PY{n}{plt}\PY{o}{.}\PY{n}{show}\PY{p}{(}\PY{p}{)}
\end{Verbatim}
\end{tcolorbox}

    \begin{center}
    \adjustimage{max size={0.9\linewidth}{0.9\paperheight}}{output_6_0.png}
    \end{center}
    { \hspace*{\fill} \\}
    
    \begin{tcolorbox}[breakable, size=fbox, boxrule=1pt, pad at break*=1mm,colback=cellbackground, colframe=cellborder]
\prompt{In}{incolor}{6}{\boxspacing}
\begin{Verbatim}[commandchars=\\\{\}]
\PY{c+c1}{\PYZsh{} Correlation Analysis}
\PY{n}{corr} \PY{o}{=} \PY{n}{np}\PY{o}{.}\PY{n}{corrcoef}\PY{p}{(}\PY{n}{velocities}\PY{p}{,} \PY{n}{distances}\PY{p}{)}\PY{p}{[}\PY{l+m+mi}{0}\PY{p}{,} \PY{l+m+mi}{1}\PY{p}{]}
\PY{n+nb}{print}\PY{p}{(}\PY{l+s+sa}{f}\PY{l+s+s2}{\PYZdq{}}\PY{l+s+s2}{Correlation between Velocity and Distance: }\PY{l+s+si}{\PYZob{}}\PY{n}{corr}\PY{l+s+si}{:}\PY{l+s+s2}{.2f}\PY{l+s+si}{\PYZcb{}}\PY{l+s+s2}{\PYZdq{}}\PY{p}{)}
\end{Verbatim}
\end{tcolorbox}

    \begin{Verbatim}[commandchars=\\\{\}]
Correlation between Velocity and Distance: 0.83
    \end{Verbatim}

    Now that we have explored the Droplet size, Velocity vs Distance and the
Impact angle, lets try and estimate the source of the spatter by reverse
engineering based on the impact angle and distances

    \begin{tcolorbox}[breakable, size=fbox, boxrule=1pt, pad at break*=1mm,colback=cellbackground, colframe=cellborder]
\prompt{In}{incolor}{9}{\boxspacing}
\begin{Verbatim}[commandchars=\\\{\}]
\PY{c+c1}{\PYZsh{} \PYZhy{}\PYZhy{}\PYZhy{} Reverse Engineering Blood Source \PYZhy{}\PYZhy{}\PYZhy{}}
\PY{n}{heights} \PY{o}{=} \PY{n}{distances} \PY{o}{*} \PY{n}{np}\PY{o}{.}\PY{n}{tan}\PY{p}{(}\PY{n}{np}\PY{o}{.}\PY{n}{radians}\PY{p}{(}\PY{n}{angles}\PY{p}{)}\PY{p}{)}

\PY{k}{def}\PY{+w}{ }\PY{n+nf}{find\PYZus{}origin}\PY{p}{(}\PY{n}{distances}\PY{p}{,} \PY{n}{angles}\PY{p}{)}\PY{p}{:}
    \PY{n}{heights} \PY{o}{=} \PY{n}{distances} \PY{o}{*} \PY{n}{np}\PY{o}{.}\PY{n}{tan}\PY{p}{(}\PY{n}{np}\PY{o}{.}\PY{n}{radians}\PY{p}{(}\PY{n}{angles}\PY{p}{)}\PY{p}{)}  \PY{c+c1}{\PYZsh{} Estimate height of impact}
    \PY{n}{avg\PYZus{}origin\PYZus{}x} \PY{o}{=} \PY{n}{np}\PY{o}{.}\PY{n}{mean}\PY{p}{(}\PY{n}{distances}\PY{p}{)}
    \PY{n}{avg\PYZus{}origin\PYZus{}y} \PY{o}{=} \PY{n}{np}\PY{o}{.}\PY{n}{mean}\PY{p}{(}\PY{n}{heights}\PY{p}{)}
    \PY{k}{return} \PY{n}{avg\PYZus{}origin\PYZus{}x}\PY{p}{,} \PY{n}{avg\PYZus{}origin\PYZus{}y}

\PY{n}{origin\PYZus{}x}\PY{p}{,} \PY{n}{origin\PYZus{}y} \PY{o}{=} \PY{n}{find\PYZus{}origin}\PY{p}{(}\PY{n}{distances}\PY{p}{,} \PY{n}{angles}\PY{p}{)}
\PY{n+nb}{print}\PY{p}{(}\PY{l+s+sa}{f}\PY{l+s+s2}{\PYZdq{}}\PY{l+s+s2}{Estimated Blood Source Location: X = }\PY{l+s+si}{\PYZob{}}\PY{n}{origin\PYZus{}x}\PY{l+s+si}{:}\PY{l+s+s2}{.2f}\PY{l+s+si}{\PYZcb{}}\PY{l+s+s2}{ cm, Y = }\PY{l+s+si}{\PYZob{}}\PY{n}{origin\PYZus{}y}\PY{l+s+si}{:}\PY{l+s+s2}{.2f}\PY{l+s+si}{\PYZcb{}}\PY{l+s+s2}{ cm}\PY{l+s+s2}{\PYZdq{}}\PY{p}{)}

\PY{c+c1}{\PYZsh{} Visualizing estimated blood source}
\PY{n}{plt}\PY{o}{.}\PY{n}{scatter}\PY{p}{(}\PY{n}{distances}\PY{p}{,} \PY{n}{heights}\PY{p}{,} \PY{n}{alpha}\PY{o}{=}\PY{l+m+mf}{0.6}\PY{p}{,} \PY{n}{label}\PY{o}{=}\PY{l+s+s1}{\PYZsq{}}\PY{l+s+s1}{Blood Drops}\PY{l+s+s1}{\PYZsq{}}\PY{p}{)}
\PY{n}{plt}\PY{o}{.}\PY{n}{scatter}\PY{p}{(}\PY{n}{origin\PYZus{}x}\PY{p}{,} \PY{n}{origin\PYZus{}y}\PY{p}{,} \PY{n}{color}\PY{o}{=}\PY{l+s+s1}{\PYZsq{}}\PY{l+s+s1}{red}\PY{l+s+s1}{\PYZsq{}}\PY{p}{,} \PY{n}{marker}\PY{o}{=}\PY{l+s+s1}{\PYZsq{}}\PY{l+s+s1}{x}\PY{l+s+s1}{\PYZsq{}}\PY{p}{,} \PY{n}{s}\PY{o}{=}\PY{l+m+mi}{100}\PY{p}{,} \PY{n}{label}\PY{o}{=}\PY{l+s+s1}{\PYZsq{}}\PY{l+s+s1}{Estimated Source}\PY{l+s+s1}{\PYZsq{}}\PY{p}{)}
\PY{n}{plt}\PY{o}{.}\PY{n}{xlabel}\PY{p}{(}\PY{l+s+s1}{\PYZsq{}}\PY{l+s+s1}{Distance (cm)}\PY{l+s+s1}{\PYZsq{}}\PY{p}{)}
\PY{n}{plt}\PY{o}{.}\PY{n}{ylabel}\PY{p}{(}\PY{l+s+s1}{\PYZsq{}}\PY{l+s+s1}{Height (cm)}\PY{l+s+s1}{\PYZsq{}}\PY{p}{)}
\PY{n}{plt}\PY{o}{.}\PY{n}{title}\PY{p}{(}\PY{l+s+s1}{\PYZsq{}}\PY{l+s+s1}{Estimated Blood Source Location}\PY{l+s+s1}{\PYZsq{}}\PY{p}{)}
\PY{n}{plt}\PY{o}{.}\PY{n}{legend}\PY{p}{(}\PY{p}{)}
\PY{n}{plt}\PY{o}{.}\PY{n}{show}\PY{p}{(}\PY{p}{)}
\end{Verbatim}
\end{tcolorbox}

    \begin{Verbatim}[commandchars=\\\{\}]
Estimated Blood Source Location: X = 10.87 cm, Y = 6.62 cm
    \end{Verbatim}

    \begin{center}
    \adjustimage{max size={0.9\linewidth}{0.9\paperheight}}{output_9_1.png}
    \end{center}
    { \hspace*{\fill} \\}
    
    \textbf{What have we done?}

\begin{enumerate}
\def\labelenumi{\arabic{enumi}.}
\tightlist
\item
  The height of impact is estimated using trigonometry;
\end{enumerate}

\[Height = Distance * \tan(Impact Angle)\]

\begin{enumerate}
\def\labelenumi{\arabic{enumi}.}
\setcounter{enumi}{1}
\tightlist
\item
  Calculated the average position of the estimated origin (X,Y)
\item
  Plotted the blood spatter and marked the estimated origin in red (X)
\end{enumerate}

    \section{Machine Learning Approach to Blood Spatter
Analysis}\label{machine-learning-approach-to-blood-spatter-analysis}

\textbf{Objective} In this section, we apply K-Means clustering, an
unsupervised machine learning technique, to analyze blood spatter
patterns. The goal is to identify natural groupings within the data
based on key characteristics such as distance from the source and impact
angle. This can help forensic investigators classify different
bloodstain patterns, potentially distinguishing between different types
of impact forces or sources.

\textbf{Why K-Means Clustering?} We chose K-Means clustering for the
following reasons:

\begin{itemize}
\tightlist
\item
  \emph{Unsupervised Learning}: Since we don't have predefined labels
  for the bloodstain data, clustering helps us discover hidden
  structures.
\item
  \emph{Efficiency}: K-Means is computationally efficient, making it
  suitable for large forensic datasets.
\item
  \emph{Interpretability}: By visualizing clusters, we can infer
  different impact dynamics and blood spatter sources.
\end{itemize}

\textbf{Methodology} \\
1. Feature Selection: We used two key features for
clustering: \\
- \emph{Distance from the source (cm)}: Helps differentiate
between high-velocity vs.~low-velocity spatters. \\
- \emph{Impact Angle
(degrees)}: Indicates the trajectory and nature of impact.\\
2. K-Means Algorithm: \\
- We set the number of clusters to three (k=3) based on expected variations in spatter patterns. The algorithm assigns each
bloodstain to the nearest cluster center. We visualize the results
using a scatter plot, where different colors represent distinct
clusters. 

    \begin{tcolorbox}[breakable, size=fbox, boxrule=1pt, pad at break*=1mm,colback=cellbackground, colframe=cellborder]
\prompt{In}{incolor}{11}{\boxspacing}
\begin{Verbatim}[commandchars=\\\{\}]
\PY{k+kn}{from}\PY{+w}{ }\PY{n+nn}{sklearn}\PY{n+nn}{.}\PY{n+nn}{cluster}\PY{+w}{ }\PY{k+kn}{import} \PY{n}{KMeans}

\PY{c+c1}{\PYZsh{} \PYZhy{}\PYZhy{}\PYZhy{} Clustering Blood Spatter Data \PYZhy{}\PYZhy{}\PYZhy{}}
\PY{n}{data\PYZus{}points} \PY{o}{=} \PY{n}{np}\PY{o}{.}\PY{n}{column\PYZus{}stack}\PY{p}{(}\PY{p}{(}\PY{n}{distances}\PY{p}{,} \PY{n}{angles}\PY{p}{)}\PY{p}{)}
\PY{n}{kmeans} \PY{o}{=} \PY{n}{KMeans}\PY{p}{(}\PY{n}{n\PYZus{}clusters}\PY{o}{=}\PY{l+m+mi}{3}\PY{p}{,} \PY{n}{random\PYZus{}state}\PY{o}{=}\PY{l+m+mi}{42}\PY{p}{)}
\PY{n}{kmeans}\PY{o}{.}\PY{n}{fit}\PY{p}{(}\PY{n}{data\PYZus{}points}\PY{p}{)}
\PY{n}{labels} \PY{o}{=} \PY{n}{kmeans}\PY{o}{.}\PY{n}{labels\PYZus{}}

\PY{c+c1}{\PYZsh{} Visualizing Clusters}
\PY{n}{plt}\PY{o}{.}\PY{n}{scatter}\PY{p}{(}\PY{n}{distances}\PY{p}{,} \PY{n}{angles}\PY{p}{,} \PY{n}{c}\PY{o}{=}\PY{n}{labels}\PY{p}{,} \PY{n}{cmap}\PY{o}{=}\PY{l+s+s1}{\PYZsq{}}\PY{l+s+s1}{viridis}\PY{l+s+s1}{\PYZsq{}}\PY{p}{,} \PY{n}{alpha}\PY{o}{=}\PY{l+m+mf}{0.7}\PY{p}{)}
\PY{n}{plt}\PY{o}{.}\PY{n}{xlabel}\PY{p}{(}\PY{l+s+s1}{\PYZsq{}}\PY{l+s+s1}{Distance (cm)}\PY{l+s+s1}{\PYZsq{}}\PY{p}{)}
\PY{n}{plt}\PY{o}{.}\PY{n}{ylabel}\PY{p}{(}\PY{l+s+s1}{\PYZsq{}}\PY{l+s+s1}{Impact Angle (Degrees)}\PY{l+s+s1}{\PYZsq{}}\PY{p}{)}
\PY{n}{plt}\PY{o}{.}\PY{n}{title}\PY{p}{(}\PY{l+s+s1}{\PYZsq{}}\PY{l+s+s1}{Clustering of Blood Spatter Data}\PY{l+s+s1}{\PYZsq{}}\PY{p}{)}
\PY{n}{plt}\PY{o}{.}\PY{n}{colorbar}\PY{p}{(}\PY{n}{label}\PY{o}{=}\PY{l+s+s1}{\PYZsq{}}\PY{l+s+s1}{Cluster Label}\PY{l+s+s1}{\PYZsq{}}\PY{p}{)}
\PY{n}{plt}\PY{o}{.}\PY{n}{show}\PY{p}{(}\PY{p}{)}
\end{Verbatim}
\end{tcolorbox}

    \begin{center}
    \adjustimage{max size={0.9\linewidth}{0.9\paperheight}}{output_12_0.png}
    \end{center}
    { \hspace*{\fill} \\}
    
    \textbf{Results \& Interpretation}

\begin{itemize}
\tightlist
\item
  The clustering process groups the blood spatters based on similar
  impact angles and distances.
\item
  Distinct clusters suggest variations in blood spatter, which could
  correspond to:
\end{itemize}

\begin{enumerate}
\def\labelenumi{\arabic{enumi}.}
\tightlist
\item
  Low-velocity impact (e.g., blunt force trauma)
\item
  Medium-velocity impact (e.g., stabbing or beating)
\item
  High-velocity impact (e.g., gunshot wounds) If needed, we can
  fine-tune the clustering by adjusting k (number of clusters) or
  incorporating additional features like droplet size or velocity.
\end{enumerate}

    \section{Assumptions in Blood Spatter
Analysis}\label{assumptions-in-blood-spatter-analysis}

Blood spatter analysis relies on several fundamental assumptions to
interpret the patterns left at a crime scene. These assumptions form the
basis for reconstructing events based on bloodstain evidence.

\begin{enumerate}
\def\labelenumi{\arabic{enumi}.}
\item
  Blood Follows Predictable Physical Laws: Blood droplets behave as
  projectiles, governed by Newtonian physics. Once a blood droplet
  leaves the source, it is primarily influenced by gravity, velocity,
  air resistance, and surface tension. The trajectory and final resting
  position of the droplet can be predicted using standard physics
  equations.
\item
  Droplet Shape and Size Correlate with Force and Impact: The size of
  blood droplets depends on the velocity of impact:
\end{enumerate}

\begin{itemize}
\tightlist
\item
  Low-velocity impact (e.g., blunt force trauma) produces larger, more
  concentrated stains.
\item
  Medium-velocity impact (e.g., beating, stabbing) results in smaller
  and more dispersed droplets.
\item
  High-velocity impact (e.g., gunshot wounds) creates fine mist-like
  spatter.
\end{itemize}

Larger stains typically indicate lower-energy events, while smaller
stains suggest higher-energy impacts.

\begin{enumerate}
\def\labelenumi{\arabic{enumi}.}
\setcounter{enumi}{2}
\item
  Impact Angle Can Be Estimated Using Geometry: Blood droplets striking
  a surface at an angle form an elliptical shape. The angle of impact
  (\(\theta\)) can be calculated using the width-to-length ratio of the
  stain: \[\theta=arcsin\left(\frac{width}{length}\right)\] A more
  elongated stain suggests a shallower impact angle, while a nearly
  circular stain indicates a perpendicular impact.
\item
  Bloodstain Clustering Represents Different Sources or Events:
  Bloodstains at a crime scene are often grouped into clusters based on
  their properties. These clusters may correspond to different:
\end{enumerate}

\begin{itemize}
\tightlist
\item
  \emph{Impact events} (e.g., multiple strikes or shots).
\item
  \emph{Sources} (e.g., victim vs.~suspect blood).
\item
  \emph{Forces applied} (e.g., gunshot vs.~blunt force trauma).
\end{itemize}

\begin{enumerate}
\def\labelenumi{\arabic{enumi}.}
\setcounter{enumi}{4}
\item
  Blood Source Estimation Assumes a Common Origin: Analysts often use
  trigonometry to estimate the point of origin of a blood spatter event.
  This assumes that all stains in a given pattern originate from one
  specific source. However, movement or multiple impacts may complicate
  the actual distribution of stains.
\item
  Surface Type Affects Spatter Behavior: The surface that blood lands on
  plays a critical role in determining stain shape and distribution.
  Smooth, non-porous surfaces (e.g., glass, tile) produce well-defined
  stains with minimal distortion. Rough or absorbent surfaces (e.g.,
  fabric, carpet) cause irregular stains due to absorption and surface
  tension effects.
\item
  Blood Patterns Can Be Reverse Engineered: By analyzing bloodstains,
  forensic experts assume they can trace back to the event that caused
  the pattern. This involves using mathematical models to estimate
  factors like height of impact, velocity, and point of origin. While
  reliable, external factors such as secondary impacts, smudging, or
  additional forces can introduce uncertainty.
\end{enumerate}

In conclusion, blood spatter analysis is a powerful forensic tool, but
it relies on fundamental assumptions about physics, biology, and
geometry. Understanding these assumptions allows forensic investigators
to interpret crime scenes accurately while also recognizing potential
limitations and uncertainties in their findings.

    \section{Conclusion}\label{conclusion}

In this report, we explored the principles of blood spatter analysis,
from understanding its forensic significance to simulating and analyzing
synthetic bloodstain data. Using machine learning techniques,
particularly K-Means clustering, we identified distinct patterns within
the data, demonstrating how computational methods can aid forensic
investigations.

Our study reaffirmed several key principles in blood spatter analysis:

\begin{itemize}
\tightlist
\item
  Blood droplet size and velocity correlate with the force applied at
  the source.
\item
  Impact angle estimation allows for the reconstruction of possible
  crime scene events.
\item
  Clustering techniques help distinguish between different spatter
  patterns, potentially separating multiple impact events or different
  sources of blood.
\end{itemize}

\subsection{Implications and
Limitations}\label{implications-and-limitations}

The application of data science and machine learning in forensic
analysis holds great promise for improving the accuracy of crime scene
reconstructions. However, there are limitations:

\begin{itemize}
\tightlist
\item
  Environmental factors (e.g., air resistance, surface properties) may
  introduce variability that is difficult to model.
\item
  Human movement and secondary transfer can complicate interpretations.
\item
  Machine learning models require large, real-world datasets to enhance
  accuracy beyond simulations.
\end{itemize}

\subsection{Future Work}\label{future-work}

To improve upon this analysis, future studies could:

\begin{itemize}
\tightlist
\item
  Use supervised learning techniques to classify bloodstain patterns
  with labeled forensic data.
\item
  Incorporate 3D reconstruction methods for more precise origin
  estimations.
\item
  Validate results with real forensic cases to assess model reliability
  in actual crime scene investigations.
\end{itemize}

By bridging forensic science with machine learning, this study
highlights the potential of computational tools in crime scene analysis.
While challenges remain, continued advancements in data-driven forensics
will undoubtedly enhance the precision and reliability of blood spatter
interpretation.


    % Add a bibliography block to the postdoc
    
    
    
\end{document}
